\documentclass[a4paper,11pt]{article}
\usepackage[utf8]{inputenc}
\usepackage[polish]{babel}
\usepackage{graphicx}
\usepackage{fancyhdr}
\usepackage{lastpage}
\usepackage{listings}
\usepackage{xcolor}
\usepackage{float}
\usepackage[margin=2.5cm]{geometry}

\lstset{
	basicstyle=\ttfamily\footnotesize,
	breaklines=true,
	frame=single,
	backgroundcolor=\color{gray!10},
	xleftmargin=5pt,
	framexleftmargin=5pt,
	numbers=left,
	numberstyle=\tiny\color{gray},
	keywordstyle=\color{blue}\bfseries,
	commentstyle=\color{green!60!black},
	stringstyle=\color{red}
}

\title{Laboratorium nr 2 -- Administracja BD\\
	\large Monitorowanie liczby rekordów w bazach danych}
\author{Paweł Myszka\\
	Nr albumu: 331720}
\date{27 października 2025}

\pagestyle{fancy}
\fancyhf{}
\fancyhead[L]{Z2 -- 331720}
\fancyhead[C]{Administracja Bazami Danych}
\fancyhead[R]{Paweł Myszka}
\fancyfoot[C]{\thepage\ / \pageref{LastPage}}

\begin{document}
	
	\maketitle
	\thispagestyle{fancy}
	
	\section{Cel laboratorium}
	
	Celem laboratorium było stworzenie systemu monitorowania liczby rekordów w tabelach użytkownika. System składa się z:
	\begin{itemize}
		\item Bazy administracyjnej przechowującej historię pomiarów
		\item Procedur automatyzujących proces zbierania danych
		\item Procedur analitycznych do wyświetlania historii i statystyk przyrostów
	\end{itemize}
	
	\section{Zadanie 1 -- Tworzenie baz danych}
	
	\begin{lstlisting}[language=SQL]
DECLARE @i INT = 1;
WHILE @i <= 5
BEGIN
DECLARE @dbname SYSNAME = N'PM_DB' + CAST(@i AS NVARCHAR(4));
IF DB_ID(@dbname) IS NULL
BEGIN
DECLARE @sql NVARCHAR(MAX) = 
N'CREATE DATABASE ' + QUOTENAME(@dbname) +
N' ON (NAME = N''' + @dbname + N'_Data'', 
FILENAME = N''D:\stud\sem 5\AdminBazDa\2\' 
+ @dbname + N'_Data.mdf'', 
SIZE = 20MB, MAXSIZE = 200MB, FILEGROWTH = 10MB)';
EXEC(@sql);
END
SET @i = @i + 1;
END
	\end{lstlisting}
	
	Utworzone bazy: \texttt{PM\_DB1}, \texttt{PM\_DB2}, \texttt{PM\_DB3}, \texttt{PM\_DB4}, \texttt{PM\_DB5}.
	
	\section{Zadanie 2 -- Utworzenie tabel w bazach}
	
	W bazach \texttt{PM\_DB1}, \texttt{PM\_DB2}, \texttt{PM\_DB3} utworzono strukturę tabel:
	
	\begin{itemize}
		\item \texttt{WOJ} -- województwa (3 rekordy)
		\item \texttt{MIASTA} -- miasta (4 rekordy)
		\item \texttt{OSOBY} -- osoby (5 rekordów bazowych)
		\item \texttt{ETATY} -- etaty (4 rekordy bazowe)
	\end{itemize}
	
	Wykorzystano procedurę \texttt{pr\_seed\_pm331720}, która dynamicznie tworzy tabele i wypełnia je danymi.

	
	\section{Zadanie 3 -- Dodanie rekordów w PM\_DB3}
	
	W bazie \texttt{PM\_DB3} dodano dodatkowe rekordy:
	
	\begin{lstlisting}[language=SQL]
USE PM_DB3;
INSERT INTO dbo.OSOBY(id_miasta, imie, nazwisko)
VALUES
((SELECT TOP 1 id_miasta FROM dbo.MIASTA 
WHERE nazwa='WARSZAWA'),'Piotr','Górski'),
((SELECT TOP 1 id_miasta FROM dbo.MIASTA 
WHERE nazwa='WARSZAWA'),'Agnieszka','Leśna');
		
-- Dodanie 3 dodatkowych etatów
DECLARE @idp INT = (SELECT TOP 1 id_osoby FROM dbo.OSOBY 
WHERE imie='Piotr' AND nazwisko='Górski');
DECLARE @ida INT = (SELECT TOP 1 id_osoby FROM dbo.OSOBY 
WHERE imie='Agnieszka' AND nazwisko='Leśna');
		
INSERT INTO dbo.ETATY(id_osoby, stanowisko, pensja, [od])
VALUES
(@idp, 'Kierownik',  11000, '20240101'),
(@idp, 'Koordynator', 8000, '20230101'),
(@ida, 'Młodszy Specjalista', 4200, '20230215');
	\end{lstlisting}
	
	Po dodaniu rekordów:
	\begin{itemize}
		\item \texttt{OSOBY}: 5 + 2 = \textbf{7 rekordów}
		\item \texttt{ETATY}: 4 + 3 = \textbf{7 rekordów}
	\end{itemize}
	
	\section{Zadanie 4 -- Śledzenie liczby rekordów}
	
	\subsection{4.1 -- Utworzenie tabel administracyjnych}
	
	Utworzono bazę \texttt{B25\_ADM\_PM331720} z dwiema tabelami:
	
	\begin{lstlisting}[language=SQL]
USE B25_ADM_PM331720;
GO
		
CREATE TABLE dbo.DB_CHECK (
check_id INT IDENTITY(1,1) PRIMARY KEY,
db_nam   NVARCHAR(100) NOT NULL,
d_stamp  DATETIME NOT NULL DEFAULT GETDATE(),
opis     NVARCHAR(200) NOT NULL,
[usr]    NVARCHAR(128) NOT NULL DEFAULT USER_NAME(),
[s_usr]  NVARCHAR(128) NOT NULL DEFAULT SUSER_NAME()
);
		
CREATE TABLE dbo.DB_CHECK_ITEMS (
check_id INT NOT NULL,
tb_nam   NVARCHAR(256) NOT NULL,
liczba_reordow INT NOT NULL,
tb_check_d_stamp DATETIME NOT NULL DEFAULT GETDATE(),
CONSTRAINT FK_DB_CHECK_ITEMS__DB_CHECK 
FOREIGN KEY (check_id) REFERENCES dbo.DB_CHECK(check_id)
);
	\end{lstlisting}
	
	\textbf{Struktura danych:}
	\begin{itemize}
		\item \texttt{DB\_CHECK} -- główna tabela z informacją o wykonaniu snapshotu (kto, kiedy, która baza)
		\item \texttt{DB\_CHECK\_ITEMS} -- szczegóły: liczba rekordów w każdej tabeli dla danego snapshotu
	\end{itemize}
	
	\subsection{4.2 -- Procedura monitorująca pojedynczą bazę}
	
	Procedura \texttt{pr\_pm\_check\_one} dla podanej bazy:
	\begin{enumerate}
		\item Wstawia rekord do \texttt{DB\_CHECK} i pobiera \texttt{check\_id}
		\item Iteruje po wszystkich tabelach użytkownika
		\item Zlicza rekordy używając \texttt{sys.partitions}
		\item Zapisuje wyniki w \texttt{DB\_CHECK\_ITEMS}
	\end{enumerate}
	
	\begin{lstlisting}
CREATE PROCEDURE dbo.pr_pm_check_one
@db sysname,
@opis NVARCHAR(200) = N'PM331720 - snapshot pojedynczej bazy'
AS
BEGIN
SET NOCOUNT ON;
		
-- Walidacja istnienia bazy
IF DB_ID(@db) IS NULL 
BEGIN 
RAISERROR(N'Baza %s nie istnieje.',16,1,@db); 
RETURN; 
END
	
-- Wstawienie rekordu z informacją o snaphocie
INSERT INTO dbo.DB_CHECK(db_nam, opis) VALUES (@db, @opis);
DECLARE @check_id INT = SCOPE_IDENTITY();
		
-- Dynamiczne SQL zliczające rekordy we wszystkich tabelach
DECLARE @sql NVARCHAR(MAX) = N'
INSERT INTO dbo.DB_CHECK_ITEMS(check_id, tb_nam, liczba_reordow)
SELECT ' + CAST(@check_id AS NVARCHAR(12)) + N', 
QUOTENAME(s.name) + ''.'' + QUOTENAME(t.name),
CAST(SUM(CASE WHEN p.index_id IN (0,1) 
THEN p.[rows] ELSE 0 END) AS INT)
FROM ' + QUOTENAME(@db) + N'.sys.tables t
JOIN ' + QUOTENAME(@db) + N'.sys.schemas s 
ON s.schema_id = t.schema_id
JOIN ' + QUOTENAME(@db) + N'.sys.partitions p 
ON p.[object_id] = t.[object_id]
WHERE t.is_ms_shipped = 0
GROUP BY s.name, t.name
ORDER BY s.name, t.name;';
		
EXEC sys.sp_executesql @sql;
END
	\end{lstlisting}
	
	\textbf{Uwaga:} Użyto \texttt{sys.partitions} zamiast \texttt{COUNT(*)} dla wydajności -- nie wymaga skanowania całej tabeli.
	
	\subsubsection{Przykład użycia}
	
	\begin{lstlisting}[language=SQL]
EXEC dbo.pr_pm_check_one @db = N'PM_DB1';
	\end{lstlisting}
	
	\subsubsection{Wyniki działania}
	
	Tabela \texttt{DB\_CHECK}:
	\begin{lstlisting}
check_id  db_nam   d_stamp                   opis
--------  -------  ------------------------  --------------------------
1         PM_DB1   2025-10-27 21:03:50.553   PM331720 - snapshot ...
	\end{lstlisting}
	
	Tabela \texttt{DB\_CHECK\_ITEMS}:
	\begin{lstlisting}
check_id  tb_nam           liczba_reordow
--------  ---------------  --------------
1         [dbo].[WOJ]      3
1         [dbo].[MIASTA]   4
1         [dbo].[OSOBY]    5
1         [dbo].[ETATY]    4
	\end{lstlisting}
	
	\subsection{4.3 -- Procedura monitorująca wszystkie bazy}
	
	Procedura \texttt{pr\_pm\_check\_all} używa kursora do iteracji po wszystkich bazach użytkownika (z wyłączeniem systemowych i administracyjnej):
	
	\begin{lstlisting}[language=SQL]
CREATE PROCEDURE dbo.pr_pm_check_all
@opis NVARCHAR(200) = N'PM331720 - snapshot wszystkich user DB'
AS
BEGIN
SET NOCOUNT ON;
DECLARE @db sysname;
		
-- Kursor po bazach użytkownika
DECLARE cur CURSOR LOCAL FAST_FORWARD FOR
SELECT name FROM sys.databases
WHERE name NOT IN (N'master', N'model', N'msdb', 
N'tempdb', N'B25_ADM_PM331720')
AND state = 0 AND is_read_only = 0 
AND name LIKE N'PM_DB%';
		
OPEN cur;
FETCH NEXT FROM cur INTO @db;
		
WHILE @@FETCH_STATUS = 0
BEGIN
BEGIN TRY
-- Wywołanie procedury dla pojedynczej bazy
EXEC dbo.pr_pm_check_one @db=@db, @opis=@opis;
PRINT N'Zliczono tabele w bazie: ' + @db;
END TRY
BEGIN CATCH
PRINT N'Pominięto bazę ' + @db + N': ' + ERROR_MESSAGE();
END CATCH;
		
FETCH NEXT FROM cur INTO @db;
END
		
CLOSE cur; 
DEALLOCATE cur;
END
	\end{lstlisting}
	
	\textbf{Logika działania:}
	\begin{enumerate}
		\item Kursor iteruje po wszystkich bazach z prefiksem \texttt{PM\_DB}
		\item Dla każdej bazy wywołuje \texttt{pr\_pm\_check\_one}
		\item Obsługa błędów (TRY-CATCH) -- jeśli baza jest niedostępna, procedura kontynuuje pracę
		\item Informuje o sukcesie/porażce dla każdej bazy
	\end{enumerate}
	
	\subsubsection{Przykład użycia}
	
	\begin{lstlisting}[language=SQL]
EXEC dbo.pr_pm_check_all;
	\end{lstlisting}
	
	\subsubsection{Wyniki}
	
	Komunikaty w konsoli:
	\begin{lstlisting}
Zliczono tabele w bazie: PM_DB1
Zliczono tabele w bazie: PM_DB2
Zliczono tabele w bazie: PM_DB3
	\end{lstlisting}
	
	Dane w \texttt{DB\_CHECK\_ITEMS} (fragment):
	\begin{lstlisting}
check_id  tb_nam           liczba_reordow
--------  ---------------  --------------
2         [dbo].[WOJ]      3
2         [dbo].[MIASTA]   4
2         [dbo].[OSOBY]    5
2         [dbo].[ETATY]    4
3         [dbo].[WOJ]      3
3         [dbo].[MIASTA]   4
3         [dbo].[OSOBY]    7
3         [dbo].[ETATY]    7
	\end{lstlisting}
	
	\subsection{4.4 -- Procedura wyświetlająca historię tabeli}
	
	Procedura \texttt{pr\_pm\_history} wyświetla wszystkie pomiary dla podanej tabeli w podanej bazie:
	
	\begin{lstlisting}[language=SQL]
CREATE PROCEDURE dbo.pr_pm_history
@db sysname,
@table NVARCHAR(256)
AS
BEGIN
SET NOCOUNT ON;
		
-- Normalizacja nazwy tabeli (usunięcie nawiasów kwadratowych)
DECLARE @t_norm NVARCHAR(256) = 
REPLACE(REPLACE(@table, N'[', N''), N']', N'');
		
-- Parsowanie schematu i nazwy tabeli
DECLARE @schema SYSNAME, @tab SYSNAME;
IF CHARINDEX('.', @t_norm) > 0
BEGIN
SET @schema = PARSENAME(@t_norm, 2);
SET @tab    = PARSENAME(@t_norm, 1);
END
ELSE
BEGIN
SET @schema = N'dbo';
SET @tab = @t_norm;
END
		
-- Wyszukanie wszystkich pomiarów z różnymi wariantami nazw
DECLARE @tb_quoted NVARCHAR(300) = 
QUOTENAME(@schema) + N'.' + QUOTENAME(@tab);
DECLARE @tab_only_q NVARCHAR(300) = QUOTENAME(@tab);
DECLARE @tb_plain NVARCHAR(300) = @schema + N'.' + @tab;
		
SELECT c.db_nam, i.tb_nam, c.check_id, 
c.d_stamp, i.liczba_reordow
FROM dbo.DB_CHECK c
JOIN dbo.DB_CHECK_ITEMS i ON i.check_id = c.check_id
WHERE c.db_nam = @db
AND (i.tb_nam = @tb_quoted
OR i.tb_nam = @tab_only_q
OR REPLACE(REPLACE(i.tb_nam, N'[', N''), N']', N'') 
IN (@tb_plain, @tab))
ORDER BY c.d_stamp, c.check_id;
		
IF @@ROWCOUNT = 0
RAISERROR(N'Brak historii dla bazy %s i tabeli %s', 
16, 1, @db, @table);
END
	\end{lstlisting}

	
	\subsubsection{Przykład użycia}
	
	\begin{lstlisting}[language=SQL]
EXEC dbo.pr_pm_history @db='PM_DB3', @table='OSOBY';
	\end{lstlisting}
	
	\subsubsection{Przykładowe wyniki}
	
	\begin{lstlisting}
db_nam  tb_nam         check_id  d_stamp                   liczba_reordow
------  -------------  --------  ------------------------  --------------
PM_DB3  [dbo].[OSOBY]  5         2025-10-27 20:15:30.120   5
PM_DB3  [dbo].[OSOBY]  8         2025-10-27 21:10:45.333   7
PM_DB3  [dbo].[OSOBY]  12        2025-10-27 22:05:12.887   7
	\end{lstlisting}
	
	Historia pokazuje, że w bazie \texttt{PM\_DB3} tabela \texttt{OSOBY} początkowo miała 5 rekordów, a po dodaniu dwóch osób (Piotr Górski, Agnieszka Leśna) ma 7 rekordów.
	
	\subsection{4.5 -- Procedura porównująca dwa snapshoty}
	
	Procedura \texttt{pr\_pm\_stats} porównuje dwa pomiary i oblicza przyrosty:
	
	\begin{lstlisting}[language=SQL]
CREATE PROCEDURE dbo.pr_pm_stats
@check_id_1 INT,
@check_id_2 INT
AS
BEGIN
SET NOCOUNT ON;
		
-- a) Walidacja check_id
IF NOT EXISTS(SELECT 1 FROM dbo.DB_CHECK 
WHERE check_id = @check_id_1)
OR NOT EXISTS(SELECT 1 FROM dbo.DB_CHECK 
WHERE check_id = @check_id_2)
BEGIN
RAISERROR(N'Check_id nie istniej w B25_ADM_PM331720.dbo.DB_CHECK',
16, 3);
RETURN;
END
		
-- b) Połączenie danych z dwóch pomiarów i obliczenie przyrostów
;WITH c1 AS (
SELECT c.db_nam, i.tb_nam, i.liczba_reordow
FROM dbo.DB_CHECK c 
JOIN dbo.DB_CHECK_ITEMS i ON i.check_id = c.check_id
WHERE c.check_id = @check_id_1
),
c2 AS (
SELECT c.db_nam, i.tb_nam, i.liczba_reordow
FROM dbo.DB_CHECK c 
JOIN dbo.DB_CHECK_ITEMS i ON i.check_id = c.check_id
WHERE c.check_id = @check_id_2
),
j AS (
SELECT COALESCE(c2.db_nam, c1.db_nam) AS baza,
COALESCE(c2.tb_nam, c1.tb_nam) AS tabela,
ISNULL(c1.liczba_reordow, 0) AS liczba_reordow_dla_check_1,
ISNULL(c2.liczba_reordow, 0) AS liczba_reordow_dla_check_2,
ISNULL(c2.liczba_reordow, 0) - ISNULL(c1.liczba_reordow, 0) 
AS przyrost
FROM c1 FULL OUTER JOIN c2
ON c1.db_nam = c2.db_nam AND c1.tb_nam = c2.tb_nam
)
-- Zwrócenie szczegółowej tabeli z przyrostami
SELECT baza, tabela, liczba_reordow_dla_check_1, 
liczba_reordow_dla_check_2, przyrost
FROM j 
ORDER BY baza, tabela;
		
;WITH mx AS (
SELECT baza, tabela, przyrost,
ROW_NUMBER() OVER (PARTITION BY baza 
ORDER BY przyrost DESC, tabela) AS rn
FROM j
)
SELECT baza, tabela AS najwiekszy_przyrost_w_bazie_pomiedzy_check1_i_check2
FROM mx 
WHERE rn = 1 
ORDER BY baza;
END
	\end{lstlisting}
	
	\textbf{Kluczowe elementy:}
	\begin{itemize}
		\item \textbf{FULL OUTER JOIN} -- obsługuje przypadki gdy tabela istniała tylko w jednym pomiarze
		\item \textbf{ISNULL} -- zabezpieczenie przed NULL przy obliczaniu przyrostu
		\item \textbf{ROW\_NUMBER()} -- wybór tabeli z największym przyrostem w każdej bazie
	\end{itemize}
	
	\subsubsection{Przykład użycia}
	
	\begin{lstlisting}[language=SQL]
-- Porównanie pomiarów przed i po dodaniu rekordów
EXEC dbo.pr_pm_stats @check_id_1=5, @check_id_2=8;
	\end{lstlisting}
	
	\subsubsection{Przykładowe wyniki}
	
	\textbf{Wynik 1 -- Szczegółowe przyrosty dla wszystkich tabel:}
	\begin{lstlisting}
baza    tabela         check_1  check_2  przyrost
------  -------------  -------  -------  --------
PM_DB3  [dbo].[ETATY]  4        7        3
PM_DB3  [dbo].[MIASTA] 4        4        0
PM_DB3  [dbo].[OSOBY]  5        7        2
PM_DB3  [dbo].[WOJ]    3        3        0
	\end{lstlisting}
	
	\textbf{Wynik 2 -- Największy przyrost w każdej bazie:}
	\begin{lstlisting}
baza    najwiekszy_przyrost_w_bazie_pomiedzy_check1_i_check2
------  -----------------------------------------------------
PM_DB3  [dbo].[ETATY]
	\end{lstlisting}
	
	Wyniki pokazują, że w bazie \texttt{PM\_DB3} największy przyrost (+3) wystąpił w tabeli \texttt{ETATY}.
	
	\subsection{4.5.1 -- Statystyki miesięczne}
	
	Procedura \texttt{ST\_M\_PM331720} wykorzystuje procedurę \texttt{pr\_pm\_stats} do analizy przyrostów w danym miesiącu:
	
	\begin{lstlisting}[language=SQL]
CREATE PROCEDURE dbo.ST_M_PM331720
@parametr_ym NCHAR(6),  -- format: YYYYMM np. '202510'
@baza NVARCHAR(100)
AS
BEGIN
SET NOCOUNT ON;
		
-- Wyszukanie najwcześniejszego pomiaru w miesiącu
DECLARE @c1 INT = (
SELECT MIN(check_id) FROM dbo.DB_CHECK 
WHERE CONVERT(NCHAR(6), d_stamp, 112) = @parametr_ym 
AND db_nam = @baza
);
		
-- Wyszukanie najpóźniejszego pomiaru w miesiącu
DECLARE @c2 INT = (
SELECT MAX(check_id) FROM dbo.DB_CHECK 
WHERE CONVERT(NCHAR(6), d_stamp, 112) = @parametr_ym 
AND db_nam = @baza
);
		
-- Walidacja
IF @c1 IS NULL OR @c2 IS NULL
BEGIN
RAISERROR(N'Brak pomiarów dla wskazanego miesiąca i bazy.', 16, 6);
RETURN;
END
		
-- Wywołanie procedury porównawczej
EXEC dbo.pr_pm_stats @check_id_1 = @c1, @check_id_2 = @c2;
END
	\end{lstlisting}
	
	\textbf{Format daty:} \texttt{CONVERT(NCHAR(6), d\_stamp, 112)} zwraca YYYYMM (np. 202510 dla października 2025)
	
	\subsubsection{Przykład użycia}
	
	\begin{lstlisting}[language=SQL]
-- Przyrosty w październiku 2025 dla PM_DB3
EXEC dbo.ST_M_PM331720 @parametr_ym=N'202510', @baza=N'PM_DB3';
	\end{lstlisting}
	
	\subsection{4.5.2 -- Statystyki dzienne}
	
	Analogiczna implementacja dla statystyk dziennych:
	
	\begin{lstlisting}[language=SQL]
CREATE PROCEDURE dbo.ST_D_PM331720
@parametr_dzien NCHAR(8),  -- format: YYYYMMDD np. '20251027'
@baza NVARCHAR(100)
AS
BEGIN
SET NOCOUNT ON;
		
-- Najwcześniejszy pomiar w dniu
DECLARE @c1 INT = (
SELECT MIN(check_id) FROM dbo.DB_CHECK 
WHERE CONVERT(NCHAR(8), d_stamp, 112) = @parametr_dzien 
AND db_nam = @baza
);
		
-- Najpóźniejszy pomiar w dniu
DECLARE @c2 INT = (
SELECT MAX(check_id) FROM dbo.DB_CHECK 
WHERE CONVERT(NCHAR(8), d_stamp, 112) = @parametr_dzien 
AND db_nam = @baza
);
		
-- Walidacja
IF @c1 IS NULL OR @c2 IS NULL
BEGIN
RAISERROR(N'Brak pomiarów dla wskazanego dnia i bazy.', 16, 7);
RETURN;
END
		
-- Wywołanie procedury porównawczej
EXEC dbo.pr_pm_stats @check_id_1 = @c1, @check_id_2 = @c2;
END
	\end{lstlisting}
	
	\textbf{Format daty:} \texttt{CONVERT(NCHAR(8), d\_stamp, 112)} zwraca YYYYMMDD (np. 20251027)
	
	\subsubsection{Przykład użycia}
	
	\begin{lstlisting}[language=SQL]
-- Przyrosty w dniu 27 października 2025 dla PM_DB3
EXEC dbo.ST_D_PM331720 @parametr_dzien=N'20251027', @baza=N'PM_DB3';
	\end{lstlisting}
	
	\section{Dodatkowe elementy implementacji}
	
	\subsection{Backup bazy administracyjnej}
	
	Automatyczny backup z timestampem w nazwie pliku:
	
	\begin{lstlisting}[language=SQL]
DECLARE @bakfile NVARCHAR(260) = 
N'D:\stud\sem 5\AdminBazDa\2\B25_ADM_PM331720_backup_' + 
REPLACE(CONVERT(NVARCHAR(20), GETDATE(), 120), ':', '-') + N'.bak';
		
BACKUP DATABASE [B25_ADM_PM331720]
TO DISK = @bakfile
WITH INIT, FORMAT, 
NAME = N'Full backup B25_ADM_PM331720 - PM331720';
	\end{lstlisting}
	
	Nazwa pliku np: \texttt{B25\_ADM\_PM331720\_backup\_2025-10-27 22-30-15.bak}
	
	\subsection{Utworzenie dedykowanego użytkownika}
	
	\begin{lstlisting}[language=SQL]
-- Login na poziomie serwera
CREATE LOGIN pm_331720_login WITH PASSWORD = N'Pm331720@2025';
		
- Użytkownik w bazie administracyjnej
USE B25_ADM_PM331720;
CREATE USER pm_331720_user FOR LOGIN pm_331720_login;
		
-- Uprawnienia: odczyt i zapis
ALTER ROLE db_datareader ADD MEMBER pm_331720_user;
ALTER ROLE db_datawriter ADD MEMBER pm_331720_user;
	\end{lstlisting}
	
	\section{Scenariusz testowy}
	
	\subsection{Krok 1: Pierwszy snapshot}
	\begin{lstlisting}[language=SQL]
EXEC dbo.pr_pm_check_all @opis=N'Snapshot początkowy';
	\end{lstlisting}
	
	Stan przed dodaniem rekordów:
	\begin{itemize}
		\item PM\_DB1: OSOBY=5, ETATY=4
		\item PM\_DB2: OSOBY=5, ETATY=4
		\item PM\_DB3: OSOBY=5, ETATY=4
	\end{itemize}
	
	\subsection{Krok 2: Dodanie rekordów w PM\_DB3}
	\begin{lstlisting}[language=SQL]
USE PM_DB3;
INSERT INTO dbo.OSOBY VALUES (...); -- +2 osoby
INSERT INTO dbo.ETATY VALUES (...); -- +3 etaty
	\end{lstlisting}
	
	\subsection{Krok 3: Drugi snapshot}
	\begin{lstlisting}[language=SQL]
EXEC dbo.pr_pm_check_all @opis=N'Po dodaniu rekordów';
	\end{lstlisting}
	
	Stan po dodaniu:
	\begin{itemize}
		\item PM\_DB1: OSOBY=5, ETATY=4 (bez zmian)
		\item PM\_DB2: OSOBY=5, ETATY=4 (bez zmian)
		\item PM\_DB3: OSOBY=7, ETATY=7 (przyrost!)
	\end{itemize}
	
	\subsection{Krok 4: Analiza przyrostów}
	\begin{lstlisting}[language=SQL]
-- Porównanie check_id=5 i check_id=8
EXEC dbo.pr_pm_stats @check_id_1=5, @check_id_2=8;
-- Historia tabeli OSOBY w PM_DB3
EXEC dbo.pr_pm_history @db='PM_DB3', @table='OSOBY';
-- Statystyki miesięczne
EXEC dbo.ST_M_PM331720 @parametr_ym=N'202510', @baza=N'PM_DB3';
	\end{lstlisting}
	
	\section{Podsumowanie}
	
	W ramach laboratorium zrealizowano kompleksowy system monitorowania liczby rekordów:

	\vspace{1cm}
	
	\begin{center}
		\textit{Paweł Myszka, nr albumu 331720}\\
		\textit{27 października 2025}
	\end{center}
	
\end{document}